\section{Alternatives and their evidence}\label{sec:alternatives}
\subsection{Mindfulness and acceptance-based treatments}
The alternative that follows from the argument is to train attention and acceptance directly rather than dispute thoughts. Brewer frames anxiety as a habit loop (trigger, worry, a feeling of control) and proposes noticing, curiosity about the reward value of worry, and a ``bigger, better offer'' \citep{brewer2021unwinding}. Evidence for the specific programme is early: in a randomised trial of an app-based programme for GAD ($n=65$), anxiety scores fell 67\% versus 14\% under usual care, with mediation implicating reduced worry \citep{roy2021app}. The trial is small, self-reported, compared with usual care rather than CBT, and has the programme's developer among its authors.

Against evidence-based comparators, mindfulness and acceptance approaches show parity, not superiority. A meta-analysis found essentially no difference ($g\approx0.00$) between mindfulness-based interventions and evidence-based treatments, often CBT \citep{goldberg2018mbi}; another found no difference from traditional CBT or behavioural therapies ($g=-.07$, nine comparisons; \citealt{khoury2013mindfulness}). For acceptance and commitment therapy (ACT) versus CBT, the pooled difference was non-significant in one meta-analysis \citep{atjak2015act}, and a sceptical reading of 60 RCTs concluded that ACT is not yet well-established for any disorder \citep{ost2014act}. In 276 adults with anxiety disorders, eight weeks of mindfulness-based stress reduction was noninferior to escitalopram, with fewer discontinuations for adverse events \citep{hoge2023mbsr}. That is parity with a first-line drug, not superiority to CBT. The convergence of camps is itself informative: if CBT, behavioural activation, ACT and mindfulness are roughly equivalent, either common factors dominate or several mechanisms lead to the same place.

\subsection{Medication in ADHD}
In a network meta-analysis of 133 double-blind trials, amphetamines (standardised effect $-0.79$) and methylphenidate ($-0.49$) outperformed placebo in adults on clinician ratings, over short follow-up \citep{cortese2018adhdmeds}. The certainty varied from high or moderate to low or very low. A Cochrane review of extended-release methylphenidate found a self-rated effect of $-0.37$, judged 20 of 21 trials at high risk of bias, and noted that all trials excluded participants with psychiatric comorbidity such as anxiety or selected previous responders \citep{boesen2022ermph}. So the pivotal medication trials say little about adults with ADHD \emph{and} anxiety.
